%%%% Preface for:
%%%% Eduardo D. Sontag,
%%%% __Mathematical Control Theory: Deterministic Finite Dimensional Systems__
%%%% Springer, NY, 1990. (396+xiii pages, ISBN 0-387-97366-4; 3-540-97366-4)
%%%% Series: Textbooks in Applied Mathematics, Number 6.
%%%% Hardcover, $39.00
%%%% Order in USA from 1-800-SPRINGER (in NJ, 201-348-4033)
%%%% save into ``book.tex'', then compile as ``latex book.tex''

\documentstyle[11pt]{article}
\topmargin -0.5in
\textheight 8.8in
\oddsidemargin 0.0in
\evensidemargin 0.0in
\textwidth 6.7in

\begin{document}

\begin{center}
{\bf Preface for}\\
{\it Mathematical Control Theory: Deterministic Finite Dimensional Systems}\\
by Eduardo D. Sontag, \copyright Springer Verlag 1990, NY\\
Series: Textbooks in Applied Mathematics, Number 6
\end{center}
This textbook introduces the basic concepts and results of mathematical
control and system theory.  Based on courses that I have taught during the
last 15 years, it presents its subject in a self-contained and elementary
fashion.  It is geared primarily to an audience consisting of mathematically
mature advanced undergraduate or beginning graduate students.  In addition, it
can be used by engineering students interested in a rigorous, proof-oriented
systems course that goes beyond the classical frequency-domain material and
more applied courses.

The minimal mathematical background that is required of the reader is a working
knowledge of linear algebra and differential equations.  Elements of the
theory of functions of a real variable, as well as elementary notions from
other areas of mathematics, are used at various points and are reviewed in the
appendixes.  However, the book was written in such a manner that readers not
comfortable with these techniques should still be able to understand all
concepts and to follow most arguments, at the cost of skipping a few technical
details and making some simplifying assumptions in a few places ---such as
dealing only with piecewise continuous functions where arbitrary measurable
functions are allowed.  In this dual mode, I have used the book in courses at
Rutgers University with mixed audiences consisting of mathematics, computer
science,
electrical engineering, and mechanical and aerospace engineering students.
Depending on the detail covered in class, it can be used for a one-, two-, or
three-semester course.  By omitting the chapter on optimal control and the
proofs of the results in the appendixes, a one-year course can be structured
with no difficulty.

The book covers what constitutes the common core of control theory: The
algebraic theory of
linear systems, including controllability, observability, feedback
equivalence, and minimality; stability via Lyapunov, as well as input/output
methods; ideas of optimal control; observers and dynamic feedback;
parameterization of stabilizing controllers (in the scalar case only); and
some very basic facts about frequency domain such as the Nyquist criterion.
Kalman filtering is introduced briefly through a deterministic version of
``optimal observation;''  this avoids having to develop the theory of
stochastic processes and represents a natural application of optimal control
and observer techniques.  In general, no stochastic or infinite dimensional
results are covered, nor is a detailed treatment given of nonlinear
differential-geometric control; for these more advanced areas, there are
many suitable texts.

The introductory chapter describes the main contents of the book in an
intuitive and informal manner, and the reader would be well advised to read
this in detail.  I suggest spending at least a week of the course covering
this material.  Sections marked with asterisks can be left out without loss of
continuity.  I have only omitted the proofs of what are labeled
``Lemma/Exercises,'' with the intention that these proofs {\it must} be worked
out
by the reader in detail.  Typically, Lemma/Exercises ask the reader to prove
that some elementary property holds, or to prove a difference equation
analogue of a differential equation result just shown.  Only a trivial to
moderate effort is required, but filling in the details forces one to
read the definitions and proofs carefully. ``Exercises'' are for the most
part also quite simple; those that are harder are marked
\def\harder{$\Diamond\;$}
with the symbol \harder.
 
Control and system theory shares with some other areas of ``modern''
applied mathematics (such as quantum field theory, geometric mechanics, and
computational
complexity), the characteristic of employing a broad range of mathematical
tools, providing many challenges and possibilities for interactions with
established areas of ``pure'' mathematics.
Papers often use techniques from algebraic and differential geometry,
functional analysis, Lie algebras, combinatorics, and other areas, in addition
to tools from complex variables, linear algebra, and ordinary and partial
differential equations.  While staying within the bounds of an introductory
presentation, I have tried to provide pointers toward some of these exciting
directions of research, particularly through the remarks at the ends of
chapters and through
references to the literature.  (At a couple of points I include
further details, such as when I discuss Lie group actions and families of
systems, or degree theory and nonlinear stabilization, or ideal theory and
finite-experiment observability, but these are restricted to small sections
that can be skipped with no loss of continuity.)
 
This book should be useful too as a research reference source, since I have
included complete proofs of various technical results often used in papers but
for which precise citations are hard to find.  I have also compiled a rather
long bibliography covering extensions of the results given and other areas
not treated, as well as a detailed index facilitating access to these.

Although there are hundreds of books dealing with various aspects of control
and system theory, including several extremely good ones, this text is unique
in its emphasis on foundational aspects.  I know of no other book that, while
covering the range of topics treated here and written in a standard
theorem/proof style, also develops the necessary techniques from scratch and
does not assume any background other than basic mathematics. In fact, much
effort was spent trying to find consistent notations and terminology across
these various areas.  On the other hand, this book does not provide realistic
engineering examples, as there are already many books that do this well.  (The
bibliography is preceded by a discussion of such other texts.)

I made an effort to highlight the distinctions as well as the similarities
between continuous- and discrete-time systems, and the process (sampling) that
is used in practice for computer control.  In this connection, I find it
highly probable that future developments in control theory will continue the
movement toward a more ``computer-oriented'' and ``digital/logical'' view of
systems, as opposed to the more classical continuous-time smooth paradigm
motivated by analogue devices. Because of this, it is imperative that a
certain minimal amount of the `abstract nonsense' of systems (abstract
definitions of systems as actions over general time sets, and so forth) be
covered: It is impossible to pose, much less solve, systems design problems
unless one first understands what a system is.  This is analogous to
understanding the definition of function as a set of ordered pairs, instead
of just thinking of functions as only those that can be expressed by explicit
formulas.

A few words about notation and numbering conventions.
Except for theorems, which are numbered consecutively, all other environments
(lemmas, definitions, and so forth) are numbered by section, while equations
are numbered by chapter.
In formal statements, I use roman characters such as $x$ to denote states at a
given instant, reserving Greek letters such as $@x$ for state trajectories,
and a similar convention is used for control values versus control functions,
and output values versus output functions;
however, in informal discussions and examples, I use roman notation for both
points and functions, the meaning being clear from the context.
\newcommand{\halmos}{\rule{1ex}{1.4ex}}
The symbol $\halmos$ marks an end of proof, while $\Box$ indicates
the end of a remark, example, etc., as well as the end of a statement whose
proof has been given earlier.

This volume represents one attempt to address two concerns raised in the
report ``Future Directions in Control Theory: A Mathematical Perspective,''
which was prepared by an international panel under the auspices of various
American funding agencies (National Science Foundation, Air Force Office of
Scientific Research, Army Research Office, and Office of Naval Research) and
was published in 1988 by the Society for Industrial and Applied Mathematics.
Two of the main recommendations of the report were that further efforts be
made to achieve conceptual unity in the field, as well as to develop better
training for students and faculty in mathematics departments interested in
being exposed to the area.  Hopefuly, this book will be useful in helping
achieve these goals.

It is hard to even begin to acknowledge all those who influenced me
professionally and from whom I learned so much, starting with Enzo Gentile and
other faculty at the Mathematics Department of the University of Buenos Aires.
Without doubt, my years at the Center for Mathematical System Theory as a
student of Rudolf Kalman were the central part of my education, and the
stimulating environment of the center was unequaled in its possibilities for
learning.  To Professor Kalman and his long-range view of the area,
originality of thought, and emphasis on critical thinking, I will always owe
major gratitude.  Others who spent considerable time at the Center, and from
whose knowledge I also benefited immensely at that time and ever since,
include Roger Brockett, Samuel Eilenberg,
Michel Fliess, Malo Hautus, Michiel Hazewinkel, Hank Hermes, Michel Heymann,
Bruce Lee, Alberto Isidori, Ed Kamen, Sanjoy Mitter, Yves Rouchaleau, Yutaka
Yamamoto, Jan Willems, and Bostwick Wyman.  In latter years I learned much
from many people, but I am especially indebted to my Rutgers colleague H\'ector
Sussmann, who introduced me to the continuous-time nonlinear theory.  While
writing this book, I received constant and very useful feedback from graduate
students at Rutgers, including Francesca Albertini, Yuandan Lin, Wen-Sheng
Liu, Guoqing Tang, Yuan Wang, and Yudi Yang.  I wish to especially thank Renee
Schwarzschild, who continuously provided me with extensive lists of misprints,
errors, and comments.  Both Pramod Khargonekar and Jack Rugh also gave me
useful suggestions at various points.

The continued and generous research support that the Air Force Office of
Scientific Research provided to me during most of my career was instrumental
in my having the possibility of really understanding much of the material
treated in this book as well as related areas.  I wish to express my sincere
gratitude for AFOSR support of such basic research.

Finally, my special thanks go to my wife Fran and to my children Laura and
David, for their patience and understanding during the long time that this
project took.

\begin{flushright}
{\it Eduardo D.\ Sontag}
\end{flushright}
\begin{flushleft}
Piscataway, NJ\\
June, 1990.
\end{flushleft}

\end{document}
